\documentclass[tikz,border=6pt]{standalone}
\usepackage[no-math]{fontspec}
\usepackage[brazil]{babel}
\usepackage{amsmath,amssymb}
\IfFileExists{C:/Windows/Fonts/Lato_Regular.ttf}{%
  \setmainfont{Lato}[
    Path = C:/Windows/Fonts/,
    UprightFont = Lato_Regular.ttf,
    BoldFont = Lato_Bold.ttf,
    ItalicFont = Lato_Italic.ttf,
    BoldItalicFont = Lato_BoldItalic.ttf
  ]
  \setsansfont{Lato}[
    Path = C:/Windows/Fonts/,
    UprightFont = Lato_Regular.ttf,
    BoldFont = Lato_Bold.ttf,
    ItalicFont = Lato_Italic.ttf,
    BoldItalicFont = Lato_BoldItalic.ttf
  ]
}{%
  \RequirePackage[default]{lato}
}
\renewcommand{\familydefault}{\sfdefault}

\usepackage{xcolor}
\usepackage{tikz}
\usepackage{pgfplots}
\pgfplotsset{compat=1.18}
\usetikzlibrary{arrows.meta,positioning,shapes.geometric,fit,backgrounds,calc,decorations.pathreplacing,patterns}

% Paleta institucional UEFS / CHF662
\definecolor{uefsNavy}{HTML}{1B2A4A}
\definecolor{uefsTeal}{HTML}{0E7C7B}
\definecolor{uefsForest}{HTML}{2D6A4F}
\definecolor{uefsGold}{HTML}{C4983B}
\definecolor{uefsGoldDark}{HTML}{8F6B1E}
\definecolor{uefsAlert}{HTML}{B02A37}
\definecolor{uefsOrange}{HTML}{D95F02}
\definecolor{uefsInk}{HTML}{22252A}
\definecolor{uefsMuted}{HTML}{5C6470}
\definecolor{uefsWater}{HTML}{2E86AB}
\definecolor{uefsSand}{HTML}{D9B382}
\definecolor{bgNavy}{HTML}{F2F5F9}
\definecolor{bgTeal}{HTML}{E9F5F4}
\definecolor{bgForest}{HTML}{EEF5F0}
\definecolor{bgGold}{HTML}{FCF7EB}
\definecolor{bgAlert}{HTML}{FDF2F2}
\definecolor{bgGray}{HTML}{F8FAFC}
\definecolor{lineGray}{HTML}{CFD8DC}

\tikzset{
  >=Stealth,
  aaCard/.style={draw=#1, fill=white, rounded corners=4pt, line width=1pt, inner sep=6pt, align=center, font=\small},
  aaBox/.style={draw=#1, rounded corners=4pt, line width=1.1pt, inner sep=6pt, align=center, font=\footnotesize\sffamily, text=uefsInk},
  aaTitle/.style={font=\normalsize\bfseries\sffamily, text=uefsNavy, align=center},
  aaSub/.style={font=\footnotesize\sffamily, text=uefsMuted, align=center},
  aaArrow/.style={->, >=Stealth, line width=1.4pt},
  aaDArrow/.style={<->, >=Stealth, line width=1.4pt}
}


\begin{document}
\begin{tikzpicture}[font=\sffamily]

% =========================================================
% PAINEL (a) Integração Físico-Biótica e Socioespacial (Venn)
% =========================================================
\node[aaTitle] at (-4.6, 4.1) {(a) Integração Físico-Biótica e Socioespacial};

% Círculos amplos e bem espaçados
\begin{scope}[shift={(-4.6, 0.2)}]
  % Meio Físico (superior esquerdo)
  \fill[uefsNavy!18, opacity=0.72] (-1.45, 0.95) circle (2.25cm);
  \draw[uefsNavy!65, line width=1pt] (-1.45, 0.95) circle (2.25cm);

  % Meio Biótico (superior direito)
  \fill[uefsForest!18, opacity=0.72] (1.45, 0.95) circle (2.25cm);
  \draw[uefsForest!75, line width=1pt] (1.45, 0.95) circle (2.25cm);

  % Sociedade & Território (inferior)
  \fill[uefsGold!22, opacity=0.72] (0, -1.25) circle (2.25cm);
  \draw[uefsGoldDark!75, line width=1pt] (0, -1.25) circle (2.25cm);

  % Rótulos bem posicionados inteiramente dentro de cada círculo
  \node[align=center, font=\footnotesize\bfseries, text=uefsNavy, text width=2.8cm]
    at (-2.15, 1.45) {Meio Físico\\[-1pt]{\normalfont\scriptsize(relevo, solo, clima, água)}};

  \node[align=center, font=\footnotesize\bfseries, text=uefsForest, text width=2.8cm]
    at (2.15, 1.45) {Meio Biótico\\[-1pt]{\normalfont\scriptsize(vegetação, fauna, biota)}};

  \node[align=center, font=\footnotesize\bfseries, text=uefsGoldDark, text width=3.4cm]
    at (0, -2.15) {Sociedade \& Território\\[-1pt]{\normalfont\scriptsize(uso da terra, técnica, política, economia)}};

  % Caixa central na interseção tríplice
  \node[draw=uefsAlert, fill=white, rounded corners=3.5pt, line width=1.2pt,
        inner sep=5pt, align=center, font=\scriptsize\bfseries, text=uefsAlert,
        text width=2.3cm]
    at (0, 0.15) {ANÁLISE\\AMBIENTAL\\INTEGRADA};
\end{scope}

% =========================================================
% PAINEL (b) Sistema Tripolar GTP (Bertrand, 2004)
% =========================================================
\node[aaTitle] at (5.3, 4.1) {(b) Sistema Tripolar GTP (Bertrand, 2004)};

\begin{scope}[shift={(5.3, 0.2)}]
  % Nó Geossistema (G)
  \node[aaBox=uefsTeal, fill=bgTeal, text width=4.0cm, minimum height=2.2cm] (G) at (-2.6, 1.45) {
    \textbf{\color{uefsTeal}GEOSSISTEMA (G)}\\[3pt]
    \scriptsize Fonte natural, dinâmica físico-ecológica e tempo da natureza
  };

  % Nó Território (T)
  \node[aaBox=uefsGold, fill=bgGold, text width=4.0cm, minimum height=2.2cm] (T) at (2.6, 1.45) {
    \textbf{\color{uefsGoldDark}TERRITÓRIO (T)}\\[3pt]
    \scriptsize Apropriação, gestão, recurso econômico e poder institucional
  };

  % Nó Paisagem (P)
  \node[aaBox=uefsNavy, fill=bgNavy, text width=4.8cm, minimum height=2.1cm] (P) at (0, -1.95) {
    \textbf{\color{uefsNavy}PAISAGEM (P)}\\[3pt]
    \scriptsize Identidade cultural, percepção social, patrimônio e leitura integrada do espaço
  };

  % Setas duplas limpas com encurtamento para não tocar as bordas
  \draw[aaDArrow, uefsNavy, shorten >=4pt, shorten <=4pt] (G.east) -- (T.west);
  \draw[aaDArrow, uefsTeal, shorten >=4pt, shorten <=4pt] ([xshift=0.35cm]G.south) -- ([xshift=-0.85cm]P.north);
  \draw[aaDArrow, uefsGoldDark, shorten >=4pt, shorten <=4pt] ([xshift=-0.35cm]T.south) -- ([xshift=0.85cm]P.north);
\end{scope}

\end{tikzpicture}
\end{document}
